\chapter{相关实验与空间观测：怎样检验模型}
\section{MRX的Hall四极场与碰撞性扫描}
\citet{Ren2005}在MRX测得面外四极磁场，并观察到Hall特征随碰撞性增加而减弱。\citet{Yamada2006}进一步讨论碰撞到弱碰撞过渡时层形态由长矩形向双楔形改变，以及重联率变化。它们提供的是对两流体分离结构的实验检验，而不是仅凭图像确定电子X点的所有非理想项。

对照模拟应提取 $B_z-B_g$、面内电流回路、离子／电子速度差以及片厚相对 $d_i$。实验有限探针间距和有限时间平均会平滑电子层，所以应对模拟做相同采样与卷积之后再比较峰值。

\section{导引场：局地Hall改变和全局背压共同作用}
\citet{GuideExp2012}系统改变外加导引场，发现四极图样和重联速率改变，局地面内Hall电流减弱以及全局导引场压缩相关背压共同参与。\citet{Dorfman2008}的几何模拟说明，上下游容积不对称与装置操作方式可改变层位置和速率。

因此实验与平板周期模拟若不同，不应首先归因于“实验有湍流”。先匹配驱动、开闭边界、下游排出条件、导引场压缩和几何体积，是更可检验的解释路径。

\section{非对称重联的能量通道}
\citet{Yoo2018}比较实验室和磁层中的非对称重联：电子能量沉积集中在电子解磁区域，离子可由面内电势结构向排气区加速。该跨尺度对照强调两流体流动与电场通道的可比性，但不会自动给出普适的电子／离子加热比例。

可用同一控制体积分比较：流入Poynting通量、流出焓和体动能、电子／离子压力应变。只有完整收支才可区分局地加热、平流输运和可逆压缩。

\section{MAST合并压缩与托卡马克相关问题}
\citet{Stanier2013}用二维双流体模拟研究MAST合并压缩启动，出现随片宽／声回旋半径变化的稳定层、开放X点和中间不稳定区。磁通绳自驱运动与局地扩散层的相互作用产生丰富现象。实验上的密度演化、合并时间及最终磁面结构提供比单一峰值率更完整的比较。

锯齿崩塌、RFP弛豫和斜撕裂的相关文献也收入索引，但真实环形系统还受磁剪切、曲率、热输运、快粒子及外部电路影响。二维Hall片中的快速磁通变化不能直接等同于托卡马克温度崩塌模型。\citep{Ottaviani1993}

\section{空间观测：Hall图样的复杂性是研究对象}
磁层卫星能测多种粒子矩和电磁场，却只沿有限轨迹穿越。\citet{Teh2026}报告导引场磁层顶相反喷流中的Hall磁场不对称，说明双极／三极／四极图样的解释需要电流重建、背景场分离和虚拟卫星检验。该文的模拟比较采用PIC，应归为观测与动理学对照，而不能计为新的Hall MHD模拟证据。

MMS事件也成为十矩模型的真实条件测试题，\citet{Closure2026}用相关不对称事件设置检验闭合对三维次级不稳定性的影响。实验／观测的最大价值是提供模型必须同时满足的多个量：密度不对称、压力张量、喷流方向、Hall场和混合，而非只要求一个速率相同。

\section{实验设计与合成诊断清单}
\begin{longtable}{P{34mm}P{58mm}P{60mm}}
\toprule 实验信号 & 流体模拟可以检验什么 & 仍需额外信息\\\midrule
四极／畸变Hall场 & 相对漂移、面内电流、背景几何 & 电子X点非理想支撑，不由图样唯一确定\\
电子与离子喷流分离 & 两种速度、不同惯性尺度 & 束流形状和非热尾需要分布函数\\
突发速率上升 & 层变薄、双稳态、磁岛触发 & 驱动历史与测量响应时间\\
电子层压力项 & 十矩分量及欧姆定律平衡 & 热流闭合验证和探针／卫星分辨率\\
三维磁通绳与湍流 & 模式谱、连接拓扑、混合层宽度 & 时间与空间变化分离、速度空间诊断\\
能量沉积 & 控制体能流、压力应变、物种内能 & 不可逆熵增及非热粒子占比\\\bottomrule
\end{longtable}
